%% ============================================================
%% 分片 24：原文第 9 节「't Hooft And Hecke Operators」（原文 glshort.tex 起于第 6755 行）
%% 由 tools/harvmac2tex.py 从 arXiv 源码（hep-th/0604151v3）自动生成。
%% 只译英文散文；公式、编号、\label、\cite 一律不改。
%% ============================================================

\subsection{The Space Of Hecke Modifications}

\label{sec:spaceheck}

The rest of this section is devoted to a closer look at the
holomorphic data in this problem.

We will first give a direct description of the holomorphic data,
emphasizing simple examples, and then explain in section
第~\ref{sec:affgrass}~节\ the more powerful framework in which the space of Hecke
modifications is usually placed in the mathematical literature.

The considerations will be local, so to keep things simple we
choose $C=\Bbb{CP}^1$, which we think of as the complex $z$-plane
plus a point at infinity.  We take the point $p_0$ to be $z=0$. To
start with, we take the gauge group to be $G=U(N)$.  We let $E_-$
be a trivial bundle of rank $n$. (Since the considerations will be
local, it does not matter what $E_-$ we start with.)  We want to
discuss how an 't Hooft/Hecke operator can modify $E_-$ to make a
new bundle $E_+$.
 We write a local holomorphic section of $E_-$
(that is, a section holomorphic near $p_0$), as
$s=(f_1,f_2,\dots,f_N)$, where the $f_i$ are local holomorphic
functions.

If $s$ is any holomorphic section of $E_-$ near $p_0$,  we can
define a line bundle ${\cal L}$ near $p_0$, as follows. ${\cal L}$
is a subbundle of $E_-$, defined, over a neighborhood of $p_0$, by
saying that sections of ${\cal L}$ are sections of $E_-$ of the
form $gs$ for an arbitrary holomorphic function $g$. We call
${\cal L}$ the line bundle generated by $s$. Now, suppose that
$s_1,\dots,s_N$ are $N$ local holomorphic sections of $E_-$ whose
restrictions to $z=0$ are linearly independent. (If we write
$s_i(z)=(f_{1i}(z),f_{2i}(z),\dots,f_{Ni}(z))$, the condition is
that the $N\times N$  matrix  $f_{ji}(0)$ has nonzero
determinant.) For each $i$, we let ${\cal L}_i$ be the line bundle
generated by $s_i$.  The linear independence ensures that near
$z=0$ we have a decomposition $E_-=\oplus_{i=1}^N{\cal L}_i$,
which is equivalent to saying that any local holomorphic section
$s$ of $E_-$ is of the form $s=\sum_i g_is_i$ with some local
holomorphic functions $g_i$.  In this situation, we say that the
sections $s_i$ generate $E_-$.

Conversely, any such local decomposition of $E_-$ takes this form.
If $E_-=\oplus_{i=1}^N{\cal L}_i$, we let $s_i$ be any section of
${\cal L}_i$ that is nonzero at $z=0$.  Then near $z=0$, ${\cal
L}_i$ is the line bundle generated by $s_i$, and so the
decomposition of $E_-$ takes the form described in the last
paragraph.

Now we consider the action of the 't Hooft/Hecke operator
$T(^L\negthinspace w)$ with $^L\negthinspace w=(m_1,m_2,\dots,
m_N)$.  With our chosen decomposition of $E$, it maps $E_-$ to the
bundle that near $z=0$ is $E_+=\oplus_i\left( {\cal L}_i\otimes
{\cal O}(p_0)^{m_i}\right)$. Away from $z=0$, $E_+$ is just the
same as $E_-$.

But ${\cal L}_i\otimes {\cal O}(p_0)^{m_i}$ is the line bundle
whose general section, near $z=0$, is $s_i/z^{m_i}$. Thus, near
$z=0$ a general section of $E_+$ takes the form $\sum_{i=1}^N
g_i(z) s_i(z)z^{-m_i}$, where the functions $g_i$ are holomorphic
near $z=0$.  Away from $z=0$, a holomorphic section of $E_+$ is
the same as a holomorphic section of $E_-$: it takes the form
$(f_1,\dots,f_N)$ with any holomorphic functions $f_1,\dots,f_N$.

Different decompositions may lead to the same $E_+$. This is
precisely what we want to investigate; we want to investigate the
space of all bundles $E_+$ to which $E_-$ can be mapped by a given
't Hooft operator.  In formulating this question, we must remember
that any Hecke modification $E_+$ of $E_-$ is naturally isomorphic
to $E_-$ away from $z=0$, so if $E_+$ and $\tilde E_+$ are two
Hecke modifications of $E_-$, they are naturally isomorphic to
each other away from $z=0$.  We consider $E_+$ and $\tilde E_+$
equivalent if and only if the natural isomorphism away from $z=0 $
can be extended over $z=0$.  We want to classify the possible
$E_+$'s up to this equivalence.  What this means in practice will
become clear as we consider examples.

We begin with a few simple examples with $N=2$. First we take
$^L\negthinspace w=(1,0)$. We suppose that
$s_1(z)=(e_1(z),f_1(z))$ and $s_2(z)=(e_2(z),f_2(z))$, with
$(e_1(0),f_1(0))$ and $(e_2(0),f_2(0))$ linearly independent.
$E_+$ is generated by $s_1'=z^{-1}s_1$ together with $s_2$.  In
particular, $s_1=zs_1'$ is a section of $E_+$, as of course is
$s_2$.  Since $s_1$ and $s_2$ generate $E_-$, this means that any
section of $E_-$ is a section of $E_+$.  But a section of $E_+$
may have a polar part proportional to $z^{-1}s_1(0)$.

In sum, in this example, any section of $E_+$ takes the form
% 式 (9.16)
\setcounter{equation}{15}
\begin{equation}\label{eq:milok}
s(z)={c\over z}s_1(0)+(h_1(z),h_2(z))
\end{equation}
 near $z=0$,
where $c$ is a complex constant and $h_1$, $h_2$ are local
holomorphic functions. Clearly, $E_+$ depends on the choice of
$s_1(0)=(e_1(0),f_1(0))$ only up to complex scaling.  So we should
think of $(e_1(0),f_1(0))$ as defining a point in $\Bbb{CP}^1$,
and this $\Bbb{CP}^1$ parametrizes all bundles $E_+$ that can be
constructed from $E_-$ by acting with $T(^L\negthinspace w)$ for
$^L\negthinspace w=(1,0)$.

We will say that $\Bbb{CP}^1$ is the space of all Hecke
modifications of type $^L\negthinspace w=(1,0)$. We denote this
space as ${{\EUY}}(^L\negthinspace w)$, or
${{\EUY}}(^L\negthinspace w,p_0)$ if we want to specify the point,
in this case $p_0$, at which the Hecke modification is made.  We
also denote the same space as ${{\EUY}}(^L\negthinspace R)$, where
$^L\negthinspace R$ is the representation of $^L\negthinspace G$
with highest weight $^L\negthinspace w$, and refer to it
alternatively as the space of Hecke modifications of type
$^L\negthinspace R$.

In this particular example, the space of Hecke modifications is
compact and smooth.  But that is not typical.  For another
example, take $^L\negthinspace w=(2,0)$.  The Hecke operator
$T(^L\negthinspace w)$, acting on $E_-$ with the same
decomposition as before, now produces a bundle $E_+$ that is
generated by $s_1'=z^{-2}s_1$ and $s_2$. Since $s_1=z^2s_1'$ is a
section of $E_+$, along with $s_2$, it is still true that any
section $(h_1,h_2)$ of $E_-$ defines a section of $E_+$.  In
addition, $E_+$ has a two-dimensional space of polar sections,
generated by $z^{-2}s_1$ and $z^{-1}s_1$.  Let us expand
$e_1(z)=u+zu'+{\cal O}(z^2)$, $f_1(z)=v+zv'+{\cal O}(z^2)$, with
$u,u',v,v'\in\Bbb{C}$. Then the general local holomorphic section
of $E_+$ takes the form % 式 (9.17)
\setcounter{equation}{16}
\begin{equation}\label{eq:ilko}
s(z)={c\over
z^2}(u+zu',v+zv')+{d\over z} (u,v)+(h_1(z),h_2(z)),
\end{equation}
 where $c$ and
$d$ are complex constants, and $h_1(z),h_2(z)$ are holomorphic
functions.

How can we describe the space of Hecke modifications in this case?
Obviously, an overall scaling of the four variables
% 式 (9.18)
\setcounter{equation}{17}
\begin{equation}\label{eq:lorky}
(u,u',v,v')\to\lambda (u,u',v,v')
\end{equation}
 is inessential; it
can be absorbed in scaling $c$ and $d$.  But there is another
equivalence relation that we should take into account. If we
transform % 式 (9.19)
\setcounter{equation}{18}
\begin{equation}\label{eq:morky}
(u',v')\to (u', v')+w(u,v),
\end{equation}
 for  $w\in
\Bbb{C}$, then this can be absorbed in $d\to d-wc$. So we should
regard this transformation as another equivalence relation on the
parameters $u,u',v,v'$.  The two symmetries both arise from the
fact that $E_+$ is invariant under $s_1\to gs_1$, where $g$ is a
holomorphic function that is invertible at $z=0$.

\def\k{b}
So the space of Hecke modifications is parametrized by
$(u,u',v,v')$, with $u$ and $v$ not both zero (since $s_1(0)$ must
be nonzero), and subject to the equivalences \eqref{eq:lorky}\ and \eqref{eq:morky}.
The invariants under the transformation \eqref{eq:morky}\ are $u$, $v$, and
$\k=vu'-uv'$.  So we can take this scaling into account by simply
using $u$, $v$, and $\k$ to parametrize the Hecke modifications.
$\k$ scales under \eqref{eq:lorky}\ with weight 2, $\k\to \lambda^2\k$. So
any triplet $(u,v,\k)$, not identically zero, defines, modulo the
scaling,  a point in a weighted projective space
$\Bbb{WCP}^2(1,1,2)$.  The space of Hecke modifications as we have
defined it so far is not all of the weighted projective space; we
are missing the single point $(0,0,1)$, since we do not allow
$u=v=0$.

Although it is possible to work with a noncompact space of Hecke
modifications, it is inconvenient.  The spaces of physical states
in the $A$-model are the cohomology groups of various moduli
spaces.  Compactness makes the definition of the appropriate
cohomology groups much more straightforward; without it, one needs
a detailed discussion of how the wavefunctions are supposed to
behave near infinity (i.e., near $(0,0,1)$ in the present case).

The space of Hecke modifications ${{\EUY}}(^L\negthinspace w)$ as
we have defined it so far does have a natural compactification,
coming from the fact that the bundle $E_+$ has a natural limit as
$u,v\to 0$. Introduce a small parameter $\epsilon$ and let
$u=\epsilon \bar u$, $v=\epsilon \bar v$, and $d=d'/\epsilon$. A
general section of $E_+$ takes the form % 式 (9.20)
\setcounter{equation}{19}
\begin{equation}\label{eq:gilko}
s(z)={c\over
z^2}(\epsilon \bar u+zu',\epsilon \bar v+zv')+{d'\over z} (\bar
u,\bar v)+(h_1(z),h_2(z)).
\end{equation}
 The limit of this as $\epsilon\to 0$,
keeping the other variables fixed, is just % 式 (9.21)
\setcounter{equation}{20}
\begin{equation}\label{eq:pilko}
s(z)={c\over z}(u',v')+{d'\over z}(\bar u,\bar
v)+(h_1(z),h_2(z)).
\end{equation}

The bundle whose general section takes this form can be given a
more simple interpretation.  The point $(0,0,1)$ on the weighted
projective space has $\k\not=0$, which is equivalent to linear
independence of the vectors $(u,v)$ and $(u',v')$ or equivalently
$(\bar u,\bar v)$ and $(u',v')$. So when we approach $(0,0,1)$ by
taking a limit as $\epsilon\to 0$, we should assume this linear
independence. That being so, \eqref{eq:pilko}\ can be described more simply
by saying that $s(z)$ can have an arbitrary simple pole at $z=0$.
Consequently, the bundle in this limit is simply $E_+=E_-\otimes
{\cal O}(p_0)$.

But this is what we get if we make a Hecke modification of the
bundle $E_-$ with $^L\negthinspace w=(1,1)$.  The Hecke operation
$T(^L\negthinspace w)$ for this $^L\negthinspace w$ is simply the
operation of tensoring with ${\cal O}(p_0)$, as we have explained
in section 第~\ref{sec:theck}~节.  The conclusion, then, is that the space of
Hecke modifications with $^L\negthinspace w=(2,0)$ has a natural
compactification in which one also allows a Hecke modification
with $^L\negthinspace w=(1,1)$.  We will write $\bar
{{\EUY}}(^L\negthinspace w)$ for this sort of compactification of
the space of Hecke modifications of type ${{\EUY}}(^L\negthinspace
w)$ (we write more explicitly $\bar {{\EUY}}(^L\negthinspace
w,p_0)$ if we wish to specify the point at which the Hecke
modification is made).

We recall that in general an 't Hooft operator is specified by the
choice of a $G$-bundle over $\Bbb{CP}^1$; once such a choice is
given, one then solves the Bogomolny equations requiring a local
singularity determined by that choice of $G$-bundle. In
particular, for $G=U(2)$, the 't Hooft operator with
$^L\negthinspace w=(m_1,m_2)$ is defined using the $U(2)$ bundle
${\cal O}(m_1)\oplus {\cal O}(m_2)$ over $\Bbb{CP}^1$.  The only
topological invariant of such a bundle is its first Chern class,
which is $m_1+m_2$. The bundle ${\cal O}(2)\oplus {\cal O}$, which
corresponds to $^L\negthinspace w=(2,0)$, is an ``unstable
bundle''; it can be infinitesimally perturbed to change its
holomorphic type to ${\cal O}(1)\oplus {\cal O}(1)$, which
corresponds to $^L\negthinspace w =(1,1)$. The analysis that we
have just described shows that the Hecke modification
corresponding to $^L\negthinspace w=(1,1)$ must be included in
order to compactify the space of Hecke modifications with
$^L\negthinspace w=(2,0)$.

The other example that we considered is quite different.  The
bundle ${\cal O}(1)\oplus {\cal O}(0)$  cannot be perturbed to any
other holomorphic type. So the space of Hecke modifications with
$^L\negthinspace w=(1,0)$  is smooth and compact.

\par\medskip{\it The Singularity of The Weighted Projective Space}

To understand more deeply the example with $^L\negthinspace
w=(2,0)$, let us reexamine the weighted projective space
$\Bbb{WCP}^2(1,1,2)$, which has homogeneous coordinates $(u,v,\k)$
of weights $(1,1,2)$. At the point $(0,0,1)$, this space has a
$\Bbb{Z}_2$ orbifold singularity.

The local structure is therefore that of the $A_1$ singularity
$\Bbb{C}^2/\Bbb{Z}_2$. $\Bbb{C}^2/\Bbb{Z}_2$ is a hyper-Kahler
orbifold.  Unlike a typical complex singularity, the $A_1$
singularity can be resolved and deformed in a compatible fashion.
This is most naturally understood  \cite{ref121} by exhibiting
$\Bbb{C}^2/\Bbb{Z}_2$ as a hyper-Kahler quotient of
$\Bbb{H}^2\cong \Bbb{C}^4$ by $U(1)$. Instead of following that
route, we will here just describe by hand the deformation and
resolution of the singularity of the weighted projective space
$\Bbb{WCP}^{2}(1,1,2)$.

To describe the deformation, we set $a_1=u^2,$ $a_2=uv$,
$a_3=v^2$. The variables $(a_1,a_2,a_3,\k)$ all have weight two,
and modulo scaling they parametrize an ordinary projective space
$\Bbb{CP}^3$.  However, they obey $a_1a_3-a_2^2=0$, so
$\Bbb{WCP}^2(1,1,2)$ is the subvariety of $\Bbb{CP}^3$ defined by
this equation.  The equation $a_1a_3-a_2^2=0$ can be deformed to
$a_1a_3-a_2^2=\epsilon \k^2$ with a small parameter $\epsilon$.
This last equation defines a smooth quadric in $\Bbb{CP}^3$, which
is isomorphic to\footnote{By a linear change of variables, one can
arrange $a_1,a_2,a_3$ and $\k$ as the matrix elements of a
$2\times 2$ matrix $M$ such that the equation becomes $\det(M)=0$.
This equation says that $M$ is of rank 1, and so takes the form
$M_{ij}=x_iy_j$, where the $x_i$ and $y_j$ are homogeneous
coordinates for, respectively, the two factors in
$\Bbb{CP}^1\times \Bbb{CP}^1$.} $\Bbb{CP}^1\times \Bbb{CP}^1$, so
$\Bbb{WCP}^2(1,1,2)$ can be deformed to $\Bbb{CP}^1\times
\Bbb{CP}^1$.

To describe the resolution, we begin with $\Bbb{C}^4$ with
coordinates $u,v,\k$, and $\k'$.  We consider the action of
$U(1)\times U(1)$  by $(u,v,\k,\k')\to (\lambda u,\lambda
v,\lambda^2 \tilde \lambda \k,\tilde\lambda \k')$, where
$|\lambda|=|\tilde\lambda|=1$.  We define the moment maps
$\mu=|u|^2+|v|^2+2|\k|^2$, $\tilde\mu= |\k|^2+|\k'|^2$, and we
take a symplectic quotient of $\Bbb{C}^4$ by $U(1)^2$.  To do
this, impose the moment map equations
% 式 (9.22)
\setcounter{equation}{21}
\begin{equation}\label{eq:ripo}
\begin{aligned}|u|^2+|v|^2+2|\k|^2 & = 1\\
                                        |\k|^2+|\k'|^2& = d\\\end{aligned}
\end{equation}

for some constant $d$, and divide by $U(1)\times U(1)$. For small
positive  $d$, the equations do not permit $u=v=0$, so the pair
$(u,v)$ can be taken as homogeneous coordinates for a copy of
$B=\Bbb{CP}^1$.  The fiber over a given point in $\Bbb{CP}^1$ is,
for small $d$, another copy of $\Bbb{CP}^1$, parametrized by $\k$
and $\k'$.  Since $(\k,\k')$ transform to $(\lambda^2\k,\k')$
under $(u,v)\to (\lambda u,\lambda v)$, $\k$ and $\k'$ take values
in the line bundles ${\cal O}(2)$ and ${\cal O}$, respectively,
over $B$. So the symplectic quotient for small positive $d$ is
fibered over $B=\Bbb{CP}^1$ with the fiber being $Y=\Bbb{P}({\cal
O}(2)\oplus {\cal O})$, that is, the projectivization of the rank
two vector bundle $\CO(2)\oplus \CO$. In particular, the
symplectic quotient for small positive $d$ is a $\Bbb{CP}^1$
bundle over $\Bbb{CP}^1$, and thus a Hirzebruch surface.

If we increase $d$, nothing happens until we reach $d=1/2$. At
that point, the Hirzebruch surface develops a singularity and
reduces to the weighted projective space $\Bbb{WCP}^2(1,1,2)$. At
$d=1/2$, we simply write the second moment map equation as
$|\k'|^2=1/2-|\k|^2$, and observe that, by virtue of the first
equation, the right hand side is nonnegative. So a solution for
$\k'$ always exists and is unique up to the
action\footnote{$\tilde\lambda$ acts freely on $\k'$ except at the
point $\k'=u=v=0$, where the action of $\tilde\lambda$ generates
the same orbit as $\lambda$.  So at $d=1/2$, upon fixing $\k'$ to
be a nonnegative real number, we can omit $\tilde\lambda$ from the
description.} of $\tilde\lambda$. The symplectic quotient at
$d=1/2$ is thus parametrized by $u,v$, and $\k$, subject to
$|u|^2+|v|^2+2|\k|^2=1$ and the equivalence $(u,v,\k)\cong
(\lambda u,\lambda v,\lambda^2 \k)$ for $|\lambda|=1$. The
combined operation gives the weighted projective space
$\Bbb{WCP}^2(1,1,2)$.

Running this in reverse, reducing $d$ from $1/2$ to a smaller
value resolves the singularity of $\Bbb{WCP}^2(1,1,2)$ to give the
Hirzebruch surface.  Moreover, this particular Hirzebruch surface
can be deformed to give a simple product $\Bbb{CP}^1\times
\Bbb{CP}^1$.  For this, we just observe that the bundle
$\CO(2)\oplus \CO$ is unstable and can be deformed to
$\CO(1)\oplus \CO(1)$.  But the projectivization of $\CO(1)\oplus
\CO(1)$ is the same as the projectivization of $\CO\oplus \CO$ and
so is a trivial $\Bbb{CP}^1$ bundle over $\Bbb{CP}^1$, that is a
product $\Bbb{CP}^1\times \Bbb{CP}^1$.

The conclusion is that, in contrast to a generic singular
algebraic variety, the singularity of $\Bbb{WCP}^2(1,1,2)$ can be
deformed or resolved to give the same result topologically, and
moreover the deformation and resolution can be carried out
simultaneously. Moreover, this is also related to the fact that
the local structure near the singularity is the hyper-Kahler $A_1$
singularity.  At the end of section 第~\ref{sec:monbub}~节\ and in section
第~\ref{sec:kahstr}~节, we will explain this behavior from the point of view of
the Bogomolny equations.
 We also in section 第~\ref{sec:opprod}~节\
will use the structure near this singularity to illustrate the
computation of the operator product expansion of 't Hooft
operators.

\par\medskip{\it General Case For $U(2)$}

The space of Hecke modifications of type $^L\negthinspace
w=(m_1,m_2)$ can be analyzed similarly.  First of all, the space
of Hecke modifications of type $(m_1+c,m_2+c)$ is the same as that
of type $(m_1,m_2)$, since adding $(c,c)$ just has the effect of
tensoring the output of the Hecke transformation with
${\CO}(p_0)^c$.  So it is enough to take $^L\negthinspace w=(a,0)$
with some integer $a$, which  we can assume nonnegative.
(Otherwise we replace $(a,0)$ by $(0,a)$ and then, adding $-a$ to
each weight, by $(-a,0)$.) In describing Hecke modifications of
type $(a,0)$, it is convenient, given a function $g$ with a pole
of order $n$ at $z=0$, to write $[g]_-$ for the polar part of $g$:
$[g]_-=g_{-n}z^{-n}+g_{-n+1}z^{-n+1}+\dots + g_{-1}z^{-1}$.

Applying a Hecke transform with $^L\negthinspace w=(a,0)$ relative
to the decomposition $E_-={\cal L}_1\oplus {\cal L}_2$ introduced
above, we get a bundle $E_+$ whose general section takes the form
% 式 (9.23)
\setcounter{equation}{22}
\begin{equation}\label{eq:indon}
s=[c_0s_1]_-+[c_1zs_1]_-+[c_2z^2s_1]_-+\dots
+[c_{a-1}z^{a-1}s_1]_-+(h_1(z),h_2(z)),
\end{equation}
 with complex coefficients
$c_0,c_1,\dots,c_{a-1}$.

The bundle $E_+$ therefore depends on $s_1$ modulo terms of order
$z^a$. This leaves a total of $2a$ complex parameters. However,
the bundle $E_+$ is unaffected if we replace $s_1$ by $gs_1$,
where $g$ is a holomorphic function that is invertible at $z=0$.
By choice of $g$, we can eliminate $a$ of the complex parameters
in $s_1$, generalizing what we did for $a=1,2$. So the space of
Hecke modifications of type $(a,0)$ has complex dimension $a$. To
compactify this space, one has to include Hecke modifications of
types $(a-1,1), \,(a-2,2),\dots, (a-[a/2],[a/2])$, which appear
when successive terms in the Taylor series expansion of $s_1$
vanish.

Some of what we have explained will be more transparent to most
physicists if we shift from $^L\negthinspace w=(a,0)$ to
$^L\negthinspace w=(a/2,-a/2)$. (The space of Hecke modifications
is of course unchanged in this shift.) As we have explained in
section 第~\ref{sec:theck}~节, weights $(a/2,-a/2)$ make sense for gauge group
$G=PSU(2)=SO(3)$, $^L\negthinspace G=SU(2)$, and moreover the 't
Hooft operator with these weights is $S$-dual to a Wilson loop of
$SU(2)$ in the representation of spin $a$.

In studying confinement in $SU(2)$ gauge theory, it is redundant
to consider Wilson loops based on an arbitrary representation $R$.
Two representations on which the center of $SU(2)$ act in the same
way are equivalent for that purpose, so the only non-trivial
representation of $SU(2)$ that one must consider is the
representation of spin $1/2$.   Dually, any 't Hooft operator of
$SO(3)$ is topologically equivalent to either the trivial one or
the operator with weights $(1/2,-1/2)$.  Precisely because the
higher 't Hooft operators are topologically unstable, they are
considered inessential in applications of 't Hooft operators (such
as 't Hooft's original work) that aim at understanding the phases
of gauge theories.

{}From our point of view, the unstable 't Hooft operators are
important, but can ``mix'' with lower operators of the same
topological type.  In section 第~\ref{sec:monbub}~节, we will explain what this
mixing means in terms of the Bogomolny equations.

The general nature of the mixing is most easily stated in case the
Lie groups $G$ and $^L\negthinspace G$ are simple.  Let
$^L\negthinspace R$ be a representation of $^L\negthinspace G$.
Let us say that a representation $^L\negthinspace R'$ of
$^L\negthinspace G$ is associated to $^L\negthinspace R$ if the
highest weight of $^L\negthinspace R'$ is a weight of
$^L\negthinspace R$ but not the highest weight. Writing
$^L\negthinspace w$ and $^L\negthinspace w'$ for the highest
weights of $^L\negthinspace R$ and $^L\negthinspace R'$, we also
say in this situation that $^L w'$ is associated to
$^L\negthinspace w$. The condition is that $^L\negthinspace
w-{}^L\negthinspace w'$ is a nonzero dominant weight of the
adjoint form of $^L\negthinspace G$.

Then in general, the natural compactification of the space
${{\EUY}}(^L\negthinspace R)$ of Hecke modifications of type
$^L\negthinspace R$ includes the Hecke modifications of type
$^L\negthinspace R'$ for all representations $^L\negthinspace R'$
associated to $^L\negthinspace R$.
 For $G=U(N)$ (or $SU(N)$ or $PSU(N)$), this
statement just means that for a weight $^L\negthinspace
w=(m_1,m_2,\dots,m_N)$, if $m_i\geq m_j+2$ for some $i,j$, the
natural compactification includes Hecke modifications with $m_i$
replaced by $m_i-1$ and $m_j$ replaced by $m_j+1$, as well as all
other weights that can be obtained by repeating this move. That
this move is needed for the compactification can be seen simply by
embedding in $U(N)$ a family of $U(2)$ Hecke modifications whereby
weights $(m_i,m_j)$ degenerate to $(m_i-1,m_j+1)$.   For more
general groups, one can argue similarly, using moves generated by
various $SU(2)$ subgroups of $^L\negthinspace G$.  The role of
these moves in the compactification can also be seen in terms of
the affine Grassmannian, a notion that we briefly introduce in
section 第~\ref{sec:affgrass}~节, or from the Bogomolny equations, as in section
第~\ref{sec:monbub}~节.

\par\medskip{\it Minuscule Representations}

Generalizing the weight $(1,0)$ of $U(2)$, we will describe a few
more examples where the space of Hecke modifications is smooth and
compact. For $G=U(N)$, let us consider as in \eqref{eq:nyto}\ the weight
$^L\negthinspace w(k)=(1,1,\dots,1,0,0,\dots,0)$, with the number
of 1's equal to $k$. Picking a decomposition
$E_-=\oplus_{i=1}^N{\cal L}_i$, where ${\cal L}_i$ is generated by
a section $s_i$, the Hecke operator $T({}^L\negthinspace w(k))$
maps $E_-$ to $E_+=\oplus_{i\leq k}\left({\cal L}_i\otimes
\CO(p_0)\right)\oplus\left(\oplus_{i>k}{\cal L}_i\right)$. The
general section of $E_+$ takes the form % 式 (9.24)
\setcounter{equation}{23}
\begin{equation}\label{eq:nson}
s={1\over
z}(a_1,\dots,a_N)+(h_1,\dots,h_N)
\end{equation}
 near $p_0$, where $\vec
a=(a_1,\dots,a_N)\in \Bbb{C}^N$ is a linear combination of
$s_1(0),\dots,s_k(0)$.  $E_+$ therefore depends on the choice of
$s_i$ precisely through the subspace $W\subset \Bbb{C}^N$
generated by $s_1(0),\dots,s_k(0)$.  $W$ can be any
$k$-dimensional subspace of $\Bbb{C}^N$, so the space
${{\EUY}}(^L\negthinspace w(k))$ of Hecke modifications of type
$^L\negthinspace w(k)$ is the Grassmannian ${\rm Gr}(k,N)$ of
complex $k$-planes in $\Bbb{C}^N$.

Clearly, then, the space of Hecke modifications of type
$^L\negthinspace w(k)$
 is smooth and compact.  This generalizes the $U(2)$
example with $^L\negthinspace w=(1,0)$.  As we noted in relation
to eqn. \eqref{eq:nyto}, $^L\negthinspace w(k)$ is the highest weight of the
representation $\wedge^k\VV$ of $^L\negthinspace G=U(N)$.  This
representation is ``minuscule,'' which means that its weights form
a single Weyl orbit.  In general, minuscule representations  are
associated to stable 't Hooft operators with smooth, compact
spaces of Hecke modifications.  Minuscule representations are
precisely those that have no associated representations, since the
highest weight of a minuscule representation is its only dominant
weight.

The highest weight $^L\negthinspace w'(k)$ of the analogous
representation $\wedge^k\VV$ of $SU(N)$ is obtained by subtracting
the constant $-k/N$ from each weight of $^L\negthinspace w(k)$.
The space of Hecke modifications is unchanged.  The
representations $\wedge^k\VV$ are minuscule and generate the
representation ring of $SU(N)$. However, $SU(N)$ is the only
simple and simply-connected Lie group whose representation ring is
generated by minuscule representations. For general simple $G$,
there is precisely one minuscule representation for each character
of the center of $G$; groups other than $SU(N)$ have relatively
small centers and few minuscule representations.

\par\medskip{\it The Dimension Of The Space Of Hecke
Modifications}

For $U(N)$ (or $SU(N)$, or $PU(N)$), let us determine the
dimension of the space of Hecke modifications of an arbitrary
type.  For weight $^L\negthinspace w=(a_1,a_2,\dots,a_N)$, with
$a_1\geq a_2\geq \dots\geq a_N$, the dimension is
% 式 (9.25)
\setcounter{equation}{24}
\begin{equation}\label{eq:ydro}
\Delta_{^L\negthinspace w}=\sum_{i<j}(a_i-a_j).
\end{equation}
 To
justify this statement, we consider $N$ sections $s_i$ that
generate the trivial rank $N$ bundle $E_-=\CO\oplus
\CO\oplus\dots\oplus\CO$ near $z=0$.  We recall that this means
that the $s_i(0)$ are linearly independent.  Then we consider the
bundle $E_+$ whose general section is
% 式 (9.26)
\setcounter{equation}{25}
\begin{equation}\label{eq:ytgo}
q=\sum_{i=1}^Nz^{-a_i}g_is_i,
\end{equation}
 where the functions
$g_i(z)$ are holomorphic at $z=0$.  The bundle $E_+$ is invariant
under % 式 (9.27)
\setcounter{equation}{26}
\begin{equation}\label{eq:gunjo}
s_i\to s_i+\sum_{j\leq
i}h_js_j+\sum_{j>i}z^{a_i-a_j}h_js_j,
\end{equation}
 for generic functions $h_j$
that are holomorphic at $z=0$. Such a transformation of the $s_i$
can be absorbed in redefining the functions $g_i$ that appear in
\eqref{eq:ytgo}. Expanding $s_i$ in a Taylor series near $z=0$, the number
of coefficients that cannot be eliminated by the transformation
\eqref{eq:gunjo}\ is $\sum_{j=i+1}^N(a_i-a_j)$.  After summing over $i$, we
arrive at the formula \eqref{eq:ydro}\ for the dimension of the space of
Hecke modifications.
